\chapter{Bẫy Công Việc Dễ}
\markboth{Bẫy Công Việc Dễ}{The Trap of Easy Work}

\section{Việc Nhẹ Dần Làm Ý Chí Mềm Đi}
\begin{truyen}{Việc Nhẹ Dần Làm Ý Chí Mềm Đi}{Easy Work Softens the Will}
\chuhoa{C}{ó người trẻ rất thích những việc không cần va chạm nhiều, không cần suy nghĩ sâu, và cũng không dễ bị chê.}
\textit{(Some young people love tasks that require little friction, little deep thought, and little exposure to criticism.)}

Những việc như vậy cho cảm giác an toàn và đỡ mệt trước mắt.
\textit{(Such work feels safe and comfortable in the short term.)}

Nhưng nếu ngày nào cũng chỉ làm phần dễ, cơ bắp nghề nghiệp sẽ yếu đi mà mình không nhận ra.
\textit{(But if every day contains only the easy part, your professional muscles weaken without you noticing.)}

Công việc quá êm trong thời gian dài có thể lấy mất sức chịu khó của một người trẻ.
\textit{(Work that stays too smooth for too long can quietly take away a young person’s capacity for effort.)}

\end{truyen}

\begin{baihoc}
Sự dễ chịu kéo dài không phải lúc nào cũng là dấu hiệu tốt. Có khi nó là dấu hiệu em đang đứng yên.
\textit{(Long comfort is not always a good sign. Sometimes it means you are standing still.)}
\end{baihoc}

\begin{ghinhoanh}
\textbf{Short English to remember:}

Long comfort is not always a good sign.

Sometimes it means you are standing still.

\end{ghinhoanh}

\begin{tuhocnhanh}
\textbf{Câu hỏi tự học / Self-study questions}

1. Vì sao việc dễ lại nguy hiểm nếu kéo dài? \textit{Why can easy work become dangerous over time?}

2. Điều gì đang yếu dần? \textit{What is getting weaker?}

3. Em có đang tránh phần khó quá lâu không? \textit{Have you been avoiding the hard part for too long?}
\end{tuhocnhanh}
\ngancach

\section{Không Bị Chê, Nhưng Cũng Không Lớn}
\begin{truyen}{Không Bị Chê, Nhưng Cũng Không Lớn}{Not Criticized, Yet Not Growing}
\chuhoa{N}{hiều người sợ bị góp ý đến mức chỉ chọn những phần việc mình chắc sẽ làm vừa đủ.}
\textit{(Many people fear criticism so much that they only choose tasks they can safely do at an average level.)}

Nhờ vậy, họ ít bị chê và ít gặp rắc rối.
\textit{(Because of that, they face fewer complaints and fewer problems.)}

Nhưng sự phát triển thật thường đi cùng những lần làm chưa tới, bị sửa, rồi làm lại tốt hơn.
\textit{(Real growth often comes with doing something imperfectly, being corrected, and improving through repetition.)}

Người né góp ý quá kỹ thường cũng né luôn con đường lớn lên.
\textit{(People who carefully avoid criticism often avoid the path of growth as well.)}

\end{truyen}

\begin{baihoc}
Ít bị chê không tự động có nghĩa là em đang tiến bộ.
\textit{(Receiving little criticism does not automatically mean you are improving.)}
\end{baihoc}

\begin{ghinhoanh}
\textbf{Short English to remember:}

Receiving little criticism does not automatically mean you are improving.

\end{ghinhoanh}

\begin{tuhocnhanh}
\textbf{Câu hỏi tự học / Self-study questions}

1. Nhân vật đang tránh điều gì? \textit{What is the character avoiding?}

2. Vì sao con đường an toàn lại ít lớn? \textit{Why does the safe path often lead to less growth?}

3. Em cần nghe thêm góp ý ở đâu? \textit{Where do you need more feedback?}
\end{tuhocnhanh}
\ngancach

\section{Ở Lâu Một Chỗ Chỉ Vì Quen Tay}
\begin{truyen}{Ở Lâu Một Chỗ Chỉ Vì Quen Tay}{Staying Too Long Just Because It Feels Familiar}
\chuhoa{C}{ó người biết công việc hiện tại không còn dạy thêm gì đáng kể, nhưng vẫn ở lại vì mọi thứ đã quen.}
\textit{(Some people know their current job no longer teaches them much, yet they stay because everything feels familiar.)}

Sự quen thuộc cho cảm giác đỡ sợ hơn những bước chuyển mới.
\textit{(Familiarity feels less frightening than a new transition.)}

Nhưng có những chỗ ở lại thêm không phải là bền, mà là trì hoãn quyết định khó.
\textit{(But sometimes staying longer is not stability. It is just delaying a hard decision.)}

Khi mình chỉ ở lại vì đã quen tay, mình đang trả tiền bằng thời gian để đổi lấy cảm giác ít lo hơn.
\textit{(When you stay only because you are used to it, you pay with time in exchange for less anxiety.)}

\end{truyen}

\begin{baihoc}
Quen không phải lúc nào cũng đúng. Có khi nó chỉ là nỗi sợ được lặp lại đủ lâu.
\textit{(Familiar does not always mean right. Sometimes it is just fear repeated for too long.)}
\end{baihoc}

\begin{ghinhoanh}
\textbf{Short English to remember:}

Familiar does not always mean right.

Sometimes it is just fear repeated for too long.

\end{ghinhoanh}

\begin{tuhocnhanh}
\textbf{Câu hỏi tự học / Self-study questions}

1. Nhân vật đang ở lại vì giá trị hay vì quen? \textit{Is the character staying for value or familiarity?}

2. Cái giá của sự ở lại này là gì? \textit{What is the cost of staying?}

3. Em có đang trì hoãn quyết định nào không? \textit{Are you delaying an important decision?}
\end{tuhocnhanh}
\ngancach

\section{Muốn Thu Nhập Ổn, Không Muốn Độ Khó Tăng}
\begin{truyen}{Muốn Thu Nhập Ổn, Không Muốn Độ Khó Tăng}{Wanting Better Income Without Greater Difficulty}
\chuhoa{C}{ó bạn trẻ mong lương khá hơn, vị trí tốt hơn, nhưng lại muốn công việc vẫn nhẹ như cũ.}
\textit{(Some young people want better income and a stronger position while still hoping the work remains just as easy.)}

Đó là một mong muốn rất người, nhưng đời thật ít chiều theo kiểu đó.
\textit{(That is a very human wish, but real life rarely works that way.)}

Giá trị cao hơn thường đi cùng áp lực cao hơn, chuẩn cao hơn, và trách nhiệm nặng hơn.
\textit{(Higher value usually comes with higher pressure, higher standards, and heavier responsibility.)}

Nếu chỉ thích phần thưởng mà ngại độ khó, ta sẽ dễ giận cuộc đời thay vì nâng mình lên cho tương xứng.
\textit{(If you want the reward but resist the difficulty, you will resent life instead of rising to match it.)}

\end{truyen}

\begin{baihoc}
Muốn bước lên một tầng mới, em phải chấp nhận trọng lượng của tầng đó.
\textit{(If you want to move to a higher level, you must accept the weight that comes with it.)}
\end{baihoc}

\begin{ghinhoanh}
\textbf{Short English to remember:}

If you want to move to a higher level, you must accept the weight that comes with it.

\end{ghinhoanh}

\begin{tuhocnhanh}
\textbf{Câu hỏi tự học / Self-study questions}

1. Nhân vật đang muốn phần thưởng nào? \textit{What reward does the character want?}

2. Trọng lượng nào đang bị né? \textit{What weight is being avoided?}

3. Em cần tăng chuẩn ở điểm nào? \textit{Where do you need to raise your standard?}
\end{tuhocnhanh}